% \newpage
\appendix
\section{Proofs of Variational Results}
\label{app:proofs}

\subsection{Proof of the ELBO (Eq.~\ref{eq:elbo})}
We provide a detailed derivation of the evidence lower bound.

\begin{proof}
Starting from the marginal likelihood (Eq.~\ref{eq:marginal_likelihood}):
\begin{align*}
\log p_\theta(x \mid \mathbf{c})
&= \log \sum_{z} p_\theta(z, x \mid \mathbf{c}) \\
&= \log \sum_{z} q_\phi(z \mid \mathbf{c}, x) \, \frac{p_\theta(z, x \mid \mathbf{c})}{q_\phi(z \mid \mathbf{c}, x)} \\
&= \log \, \mathbb{E}_{q_\phi(z \mid \mathbf{c}, x)} \!\left[ \frac{p_\theta(z, x \mid \mathbf{c})}{q_\phi(z \mid \mathbf{c}, x)} \right].
\end{align*}

Applying Jensen’s inequality yields
\begin{align*}
\log p_\theta(x \mid \mathbf{c})
&\geq \mathbb{E}_{q_\phi(z \mid \mathbf{c}, x)} \!\left[ \log \frac{p_\theta(z, x \mid \mathbf{c})}{q_\phi(z \mid \mathbf{c}, x)} \right] \\
&= \mathbb{E}_{q_\phi(z \mid \mathbf{c}, x)} \!\big[\log p_\theta(z, x \mid \mathbf{c})\big] 
- \mathbb{E}_{q_\phi(z \mid \mathbf{c}, x)} \!\big[\log q_\phi(z \mid \mathbf{c}, x)\big].
\end{align*}

Adding and subtracting $\mathbb{E}_{q_\phi}\big[\log p_\theta(z \mid \mathbf{c})\big]$ gives
\begin{align*}
\log p_\theta(x \mid \mathbf{c})
&\geq \mathbb{E}_{q_\phi(z \mid \mathbf{c}, x)} \!\big[\log p_\theta(z, x \mid \mathbf{c})\big] 
- \mathrm{KL}\!\left(q_\phi(z \mid \mathbf{c}, x)\,\|\,p_\theta(z \mid \mathbf{c})\right),
\end{align*}
which is precisely Eq.~\ref{eq:elbo}.
\end{proof}

\subsection{Proof of the Optimal Posterior (Eq.~\ref{eq:posterior})}
We now establish the form of the optimal approximate posterior.

\begin{proof}
Rewriting Eq.~\ref{eq:elbo}, we obtain
\[
\log p_\theta(x \mid \mathbf{c})
= \mathcal{L}(q_\phi) 
+ \mathrm{KL}\!\left(q_\phi(z \mid \mathbf{c}, x)\,\|\,p_\theta(z \mid \mathbf{c}, x)\right),
\]
where $\mathcal{L}(q_\phi)$ is the ELBO. Since the KL divergence is non-negative, the ELBO is maximized when
\[
q_\phi^\star(z \mid \mathbf{c}, x) = p_\theta(z \mid \mathbf{c}, x).
\]
By Bayes’ rule,
\[
p_\theta(z \mid \mathbf{c}, x) = \frac{p_\theta(z, x \mid \mathbf{c})}{p_\theta(x \mid \mathbf{c})}
\;\propto\; p_\theta(z, x \mid \mathbf{c})
= p_\theta(x \mid z, \mathbf{c}) \, p_\theta(z \mid \mathbf{c}),
\]
which matches Eq.~\ref{eq:posterior}.
\end{proof}


\section{Additional Results on Ring-Lite}
\label{sec:ringlite}

We further assess generalization under the self-play setting, where training is conducted via iterative SFT. We initialize from \texttt{Ring-Lite-2506}~\citep{team2025ring}, a 16.8B-parameter model with 2.75B activated parameters, and compare its performance with our framework on mathematics (AIME~24/25) and programming (LiveCodeBench v5; 2408–2502). As shown in Table~\ref{tab:ringlite}, PromptCoT~2.0 attains consistent improvements over \texttt{Ring-Lite-2506} across all benchmarks (+3.2 on AIME~24, +1.3 on AIME~25, and +2.0 on LiveCodeBench v5), indicating that the EM-driven rationale–prompt co-optimization produces synthesized problems that transfer effectively within the self-play regime.

\begin{table}[h]
\centering
\small
\begin{tabular}{lccc}
\toprule
\textbf{Model} & \textbf{AIME 24} & \textbf{AIME 25} & \textbf{LCB v5 (2408–2502)} \\
\midrule
\textbf{Ring-Lite-2506} & 76.6 & 69.1 & 60.7 \\
\textbf{PromptCoT~2.0} & \textbf{79.8} & \textbf{70.4} & \textbf{62.7} \\
\bottomrule
\end{tabular}
\caption{Evaluation results on AIME~24/25 and LiveCodeBench v5 (2408-2502) under the Self-Play setting. Bold values denote the better performance for each benchmark.}
\label{tab:ringlite}
\end{table}



\section{Instruction for Concept Extraction}
\label{app:concept_prompt}


As part of the cold-start stage, this prompt is designed to extract domain-relevant concepts from each seed problem. It instructs the language model to identify the key mathematical or programming concepts that underlie the given task.

\begin{tcolorbox}[colback=blue!5!white, colframe=blue!75!black, title=Concept Extraction Prompt]
As an expert in educational assessment, analyze this problem:

\{problem\}

Break down and identify \{num\_concepts\} foundational concepts being tested. List these knowledge points that:
\begin{itemize}
    \item Are core curriculum concepts typically taught in standard courses,
    \item Are precise and measurable (not vague like "understanding math"),
    \item Are essential building blocks needed to solve this problem,
    \item Represent fundamental principles rather than problem-specific techniques.
\end{itemize}

Think through your analysis step by step, then format your response as a Python code snippet containing a list of \{num\_concepts\} strings, where each string clearly describes one fundamental knowledge point.
\end{tcolorbox}


\section{Instruction for Rationale Generation}
\label{app:rationale_prompt}

In the cold-start stage, the following prompt instructs the language model to produce a detailed reasoning process that serves as the rationale for prompt synthesis. 

\begin{tcolorbox}[colback=blue!5!white, colframe=blue!75!black, title=Rationale Generation Instruction]
Imagine you are an expert in educational problem design.

You will be shown these components:

\textbf{Problem:} \{problem\}

\textbf{Fundamental Concepts:} \{list of concepts\}

\textbf{Difficulty Level:} \{difficulty\_level\}

Your task is to reverse-engineer a clear thinking process that shows how a teacher might design this problem. This thinking process should:
\begin{itemize}
    \item Show how combining the given foundational concepts naturally leads to a problem at the specified difficulty level.
    \item Include all key decisions and reasoning that shaped the problem design.
    \item \textbf{(IMPORTANT)} Be so precise and detailed that another teacher following these exact steps would recreate the identical problem.
    \item \textbf{(IMPORTANT)} Be so natural and logical that another teacher could derive the same thinking process using only the foundational concepts and difficulty level.
\end{itemize}

Present your answer after ``Thinking Process:'' with the complete step-by-step thinking process described above.
\end{tcolorbox}


\section{Case Study}




\begin{figure}[ht]
\small
    \centering
    \fbox{%
        \begin{minipage}{0.99\textwidth}
            \vspace{0.5em}
            \textbf{Foundational Concepts:}
            \begin{enumerate}
                \item Understanding the definition and properties of divisors of a number
                \item Comprehending the concept of a function and its application to number theory, specifically the sum of divisors function $f(n)$
                \item Ability to compute the sum of a sequence of numbers, including understanding the concept of cumulative sums as represented by $F(n)$
                \item Knowledge of the parity (odd or even) of numbers and how it changes under addition, particularly in the context of sums of divisors
                \item Application of modular arithmetic to determine conditions under which $F(n)$ is odd, especially modulo 2 operations
            \end{enumerate}
            
            \vspace{1em}
            \textbf{Rationale:}
            
            \noindent\textbf{Step 1: Identify the Foundational Concepts} \\
            The problem uses divisors and the sum of divisors function (Concepts 1 and 2), cumulative sums (Concept 3), parity (Concept 4), and modular arithmetic (Concept 5).
            
            \vspace{0.5em}
            \noindent\textbf{Step 2: Difficulty Level} \\
            The problem is intended for the AIME, requiring multi-step reasoning in number theory.

            \vspace{0.5em}
            \noindent\textbf{Step 3: Formulation} \\
            Define $f(n)$ as the sum of divisors of $n$. Define $F(n)=\sum_{k=1}^n f(k)$. We want the count of $n\leq 2005$ for which $F(n)$ is odd.

            \vspace{0.5em}
            \noindent\textbf{Step 4: Analysis of Parity} \\
            The parity of $F(n)$ is determined by the parity of $\sum_{k=1}^n f(k)$. Thus, we must characterize when $f(k)$ is odd.

            \vspace{0.5em}
            \noindent\textbf{Step 5: Characterizing Odd $f(k)$} \\
            A number $k$ has an odd sum of divisors exactly when $k$ is a perfect square or twice a perfect square.  
            \begin{itemize}
                \item This follows from the multiplicativity of $\sigma(n)$ and the fact that for odd primes $p$, $\sigma(p^e)$ is odd iff $e$ is even.  
                \item Therefore, the odd part of $k$ must be a square, and the power of $2$ can be $0$ or $1$.  
            \end{itemize}

            \vspace{0.5em}
            \noindent\textbf{Step 6: Counting Relevant $k$} \\
            Hence
            \[
            F(n) \equiv \#\{k\leq n:\ k\text{ is a square}\}+\#\{k\leq n:\ k=2m^2\}\pmod 2.
            \]
            Equivalently,
            \[
            F(n)\equiv \lfloor \sqrt{n}\rfloor + \lfloor \sqrt{n/2}\rfloor \pmod 2.
            \]
            So $F(n)$ is odd exactly when $\lfloor \sqrt{n}\rfloor$ and $\lfloor \sqrt{n/2}\rfloor$ have opposite parity.

            \vspace{0.5em}
            \noindent\textbf{Step 7: Solving the Problem} \\
            For $n\leq 2005$:  
            \begin{itemize}
                \item $\lfloor \sqrt{2005}\rfloor = 44$  
                \item $\lfloor \sqrt{2005/2}\rfloor = 31$
            \end{itemize}
            Counting by intervals $n\in[t^2,(t+1)^2-1]$ for $t=1,\dots,44$, and checking where $\lfloor \sqrt{n/2}\rfloor$ has opposite parity to $t$, gives the total number of integers $n\leq 2005$ with $F(n)$ odd:
            \[
            \#\{n\leq 2005 : F(n)\text{ odd}\} = 1018.
            \]

            \vspace{1em}
            \textbf{Final Problem:} \\
            For any positive integer $n$, let $f(n)$ be the sum of the divisors of $n$, and let $F(n)$ be the sum of the numbers $f(k)$ for $k = 1$ to $n$. For example, $F(4) = f(1) + f(2) + f(3) + f(4) = 1 + 3 + 4 + 7 = 15$. Find the number of positive integers $n \leq 2005$ with $F(n)$ odd.
        \end{minipage}%
    }
    \caption{Case study of rationale and problem generation for a number theory AIME-level problem.}
    \label{fig:case_study_rationale1}
\end{figure}



\begin{figure}[htb]
\small
    \centering
    \fbox{%
        \begin{minipage}{0.99\textwidth}
            \vspace{0.5em}
            \textbf{Foundational Concepts:}
            \begin{enumerate}
                \item Understanding of the properties of a cube, including its edges, vertices, and faces
                \item Knowledge of how a plane can intersect a three-dimensional solid, specifically a cube, and the types of polygons that can be formed
                \item Ability to visualize and analyze geometric shapes in three dimensions, particularly the cross-sections of a cube
                \item Understanding of the concept of perimeter and how to calculate it for different types of polygons, especially in the context of a plane intersecting a cube
                \item Application of the Pythagorean theorem to calculate distances within the cube, which are essential for determining the perimeter of the intersection polygon
            \end{enumerate}
            
            \vspace{1em}
            \textbf{Rationale:}
            
            \noindent\textbf{Step 1: Core Concept and Difficulty Level} \\
            The problem focuses on the intersection of a plane with a cube and the polygon formed by this intersection. The goal is to determine the maximum possible perimeter of such a polygon. This is appropriate for the AIME level, which requires multi-step reasoning in three dimensions.

            \vspace{0.5em}
            \noindent\textbf{Step 2: Cube and Plane Setup} \\
            We consider a cube of edge length 1. The intersecting plane is restricted from passing through any cube vertices, ensuring the cross-section is a polygon with between 3 and 6 sides.

            \vspace{0.5em}
            \noindent\textbf{Step 3: Possible Intersections} \\
            A plane intersecting a cube can produce polygons with 3, 4, 5, or 6 sides. To maximize perimeter, we must look for the configuration that yields the longest polygonal boundary.

            \vspace{0.5em}
            \noindent\textbf{Step 4: Quadrilateral Case} \\
            If the intersection is a quadrilateral, the largest perimeter arises when the cross-section is a square taken through the center of the cube. In this case each side has length 1, giving perimeter:
            \[
            P = 4.
            \]

            \vspace{0.5em}
            \noindent\textbf{Step 5: Hexagon Case} \\
            When the plane is perpendicular to a space diagonal of the cube and passes through its center, the intersection is a regular hexagon. Each side of this hexagon has length $\tfrac{\sqrt{2}}{2}$, so the perimeter is:
            \[
            P = 6 \cdot \tfrac{\sqrt{2}}{2} = 3\sqrt{2}.
            \]

            \vspace{0.5em}
            \noindent\textbf{Step 6: Conclusion} \\
            Comparing these cases, the hexagon provides the largest possible perimeter. Therefore:
            \[
            \text{Maximum Perimeter} = 3\sqrt{2}.
            \]

            \vspace{1em}
            \textbf{Final Problem:} \\
            A plane intersects a cube with edge length 1. The plane does not pass through any of the cube's vertices. What is the maximum possible perimeter of the polygon formed by the intersection of the plane and the cube?
        \end{minipage}%
    }
    \caption{Case study of rationale and problem generation for a 3D geometry AIME-level problem.}
    \label{fig:case_study_rationale2}
\end{figure}


\begin{figure}[htb]
\small
    \centering
    \fbox{%
        \begin{minipage}{0.99\textwidth}
            \vspace{0.5em}
            \textbf{Foundational Concepts:}
            \begin{enumerate}
                \item Understanding of basic arithmetic operations, specifically addition, and the concept of carrying in multi-digit addition
                \item Ability to manipulate and process arrays, including iterating through elements and performing operations on pairs of elements
                \item Comprehension of nested loops and their use in computing pairwise operations over a range of elements in an array
                \item Knowledge of modular arithmetic and digit manipulation to handle addition without carrying, focusing on individual digits
                \item Understanding of summation notation and the ability to compute the sum of a series of values, particularly in the context of nested summations
            \end{enumerate}
            
            \vspace{1em}
            \textbf{Rationale:}
            
            \noindent\textbf{Step 1: Identify the Foundational Concepts} \\
            The problem directly draws from addition without carrying (Concept 1), array iteration (Concept 2), nested loops (Concept 3), digit-wise modular arithmetic (Concept 4), and summation notation (Concept 5).

            \vspace{0.5em}
            \noindent\textbf{Step 2: Define the Problem Context} \\
            A student (Tot) adds numbers digit by digit without carrying. The task is to generalize this operation to arrays and compute the total over all pairs.

            \vspace{0.5em}
            \noindent\textbf{Step 3: Problem Statement} \\
            Given an array $a$ of length $n$, compute $\sum_{i=1}^{n}\sum_{j=i}^{n} f(a_i,a_j)$ where $f(x,y)$ is the digit-wise sum of $x$ and $y$ without carrying.

            \vspace{0.5em}
            \noindent\textbf{Step 4: Input and Output} \\
            The input consists of multiple queries. Each query specifies $n$ and then $n$ integers. For each query, output the required double summation.

            \vspace{0.5em}
            \noindent\textbf{Step 5: Constraints} \\
            The total $n$ across all queries is at most $10^5$, ensuring feasibility with efficient implementations.

            \vspace{0.5em}
            \noindent\textbf{Step 6: Sample Input/Output} \\
            Verified with examples:
            \begin{itemize}
                \item $f(53,19)=612$, matching digit-wise no-carry addition.
                \item Summations over arrays give large outputs consistent with constraints.
            \end{itemize}

            \vspace{0.5em}
            \noindent\textbf{Step 7: Difficulty Level} \\
            At the Codeforces level, the challenge is balancing array operations and efficiency. A naive $O(n^2)$ summation may be too slow, so solutions must exploit structure.

            \vspace{0.5em}
            \noindent\textbf{Step 8: Logical Flow} \\
            The reasoning starts from digit-wise addition, extends to arrays, introduces nested summations, and ensures scalability under the constraints.

            \vspace{0.5em}
            \noindent\textbf{Step 9: Final Problem Statement} \\
            \textit{Tot was doing homework in arithmetic. He was taught how to add two numbers, but not how to carry. For example, for $53+19$, he would do:}
            \[
            \begin{array}{r}
            53 \\
            +\,19 \\
            \hline
            612
            \end{array}
            \]
            \textit{Tot was given an array $a$ of length $n$. He was asked to find $\sum_{i=1}^{n}\sum_{j=i}^{n} f(a_i,a_j)$, where $f(x,y)$ is the number obtained by adding $x$ and $y$ without carrying. For example, $f(53,19)=612$.}

            \vspace{0.5em}
            \textbf{Input:} \\
            The first line contains one integer $t$ ($1 \leq t \leq 1000$) — the number of queries. Each query begins with an integer $n$ ($1 \leq n \leq 10^5$). The next line contains $n$ integers $a_1,\dots,a_n$ ($1 \leq a_i \leq 10^9$). The total $n$ over all queries does not exceed $10^5$.

            \vspace{0.5em}
            \textbf{Output:} \\
            For each query, print the required sum.

        \end{minipage}%
    }
    \caption{Case study of rationale and problem generation for a Codeforces-level programming problem.}
    \label{fig:case_study_rationale3}
\end{figure}



